\subsection{Data Processing and Visualization Tools}
\label{subsec:data_processing_and_visualization_tools}

This adaptive approach to anomaly detection ensures that detection mechanisms remain robust and reliable, even as EVSE infrastructures and their operational conditions continue to evolve. This chapter introduces the data processing and visualization tools used to implement such adaptable anomaly detection strategies.\\
% Pandas
To convert the data into more useful formats, \textit{Pandas} was used, a fast, powerful, and easy-to-use open-source data analysis tool for Python. It provides expressive data structures such as DataFrame and Series for working with structured and time series data \cite{pandas_wiki}.\\
\texttt{pandas} was primarily used to convert the data, which was originally stored in \texttt{CSV} files, into \texttt{PARQUET} files. Additionally, it allowed manipulation of the dataframes by dropping cells, adding information such as attack or benign tags to almost 7000 rows. This preprocessing step was essential to clean the data before applying machine learning algorithms \cite{pandas_wiki}.\\

% Matplotlib
To create diagrams that help better understand the data, a popular plotting library in Python called \textit{Matplotlib} was used. Specifically, the \texttt{matplotlib} module was employed, which is a state-based interface to the \texttt{matplotlib} library and provides a MATLAB-like way of creating static, animated, and interactive visualizations in Python.\\
\texttt{matplotlib} is well-suited for quickly prototyping visualizations thanks to its simple and readable syntax. This allowed the creation of a wide range of plots, including Line Plots, Scatter Plots, Bar Charts, Histograms, and Heatmaps.\\
These visualizations enabled a better understanding of the data structure, identification of patterns, and preparation of the dataset for further processing.\\

% Seaborn
In addition to \texttt{matplotlib} and \texttt{pandas}, the open-source Python library \textit{Seaborn} was also used. Seaborn is especially known for creating statistical graphics and is built on top of the \texttt{matplotlib} library. It provides a high-level interface for drawing informative and attractive statistical plots.\\
Seaborn is closely integrated with \texttt{pandas} data structures, which facilitates the visualization of entire datasets. This integration enables the creation of complex visualizations with minimal code. Key features of Seaborn include:

\begin{itemize}
    \item Plotting of complex relationships between multiple variables.
    \item Automatic aggregation and statistical estimation (e.g., means, confidence intervals).
    \item Thematic styling and enhanced color palettes for better aesthetics.
\end{itemize}
For these reasons, Seaborn was a perfect addition to the \texttt{pandas} and \texttt{matplotlib} libraries in the data analysis workflow \cite{seaborn_wiki}.